\documentclass[11pt]{article}
\usepackage{varnernotes}

\newcommand{\B}{\mathcal{B}}

\begin{document}

\lecturenotes{2b}{Is Risk-Free Really Risk-Free? Malkiel's Bond Pricing Theorems}{Fall 2026}{September 3, 2026}{Jeffrey D. Varner}

\learningobjectives
  {Prove price sensitivity}
  {Differentiate a fixed-income cash-flow model to establish the signs of its yield sensitivity and curvature.}
  {Compute local risk measures}
  {Calculate modified duration and convexity from a bond's discounted cash flows and use them to approximate a price change.}
  {Compare bond structures}
  {Predict how maturity and coupon placement affect the interest-rate exposure of otherwise comparable securities.}

\section{A Safe Payment Can Have a Risky Price}

L2a treated promised Treasury payments as known, but a holder who sells before maturity receives a market price. That price changes whenever the return available on comparable new securities changes. A fixed promise and a fluctuating resale value are therefore entirely compatible.

The 2023 failure of Silicon Valley Bank provides the motivating systems lesson. The bank combined long-duration fixed-rate assets, concentrated uninsured funding, and inadequate interest-rate and liquidity risk management. Rising rates reduced the fair value of its securities; deposit withdrawals then forced the institution to confront losses that a hold-to-maturity story had obscured. Bond sensitivity is thus one component of a coupled balance-sheet and liquidity problem, not a complete explanation by itself.

Burton Malkiel's bond-pricing principles organize the central comparative statics: price moves opposite yield, the price-yield curve is convex, longer maturity generally increases sensitivity, and a larger coupon generally reduces percentage sensitivity. These statements require an option-free fixed cash-flow model and comparisons that hold the other contract terms fixed.

\section{Price as a Function of Yield}

To isolate interest-rate risk, freeze the promised cash flows and vary only the yield used to discount them. The resulting yield-based valuation is defined locally in \defnref{yield-price}.

\begin{definition}{Yield-based bond price}{yield-price}
Let the positive integer $n$ denote the number of coupon intervals per year, let the positive integer $N$ denote the number of remaining payment dates, and let $X_j>0$ dollars be the fixed cash flow paid at coupon index $j\in\{1,\ldots,N\}$. Let the nominal annual yield to maturity be $r$ with $1+r/n>0$. Under a flat-yield convention, the time-0 dirty price $\B(r)$---the value including any accrued interest---is given by \eqnref{yield-price}:
\begin{equation}
  \label{eq:yield-price}
  \B(r):=\sum_{j=1}^{N}X_j\left(1+\frac{r}{n}\right)^{-j}.
\end{equation}
The yield $r$ is the single internal discount rate applied to every promised payment; it is not a complete term structure of spot rates.
\end{definition}

The common factor $q:=1+r/n$ converts the notation in \eqnref{yield-price} into a sum of positive discounted claims. Differentiating that sum establishes both the direction and curvature of the price response (\thmref{yield-shape}).

\begin{theorem}{Monotonicity and convexity of an option-free bond}{yield-shape}
Under the assumptions of \defnref{yield-price}, let $q:=1+r/n>0$. The first and second derivatives of the dirty price with respect to the nominal annual yield are given by \eqnref{yield-derivatives}:
\begin{equation}
  \label{eq:yield-derivatives}
  \frac{d\B}{dr}
  =-\frac{1}{n}\sum_{j=1}^{N}jX_jq^{-j-1}<0,
  \qquad
  \frac{d^2\B}{dr^2}
  =\frac{1}{n^2}\sum_{j=1}^{N}j(j+1)X_jq^{-j-2}>0.
\end{equation}
Consequently, the bond price is strictly decreasing and strictly convex as a function of yield. The result assumes fixed positive cash flows and therefore does not automatically apply to securities whose cash flows change when rates change, such as callable bonds.
\end{theorem}

The signs in \eqnref{yield-derivatives} explain the price-yield geometry: a higher yield lowers every discounted payment, while positive curvature makes the gain from a yield decrease larger than the loss from an equal yield increase, measured from the same starting yield (\figref{price-yield}).

\begin{figure}[ht]
  \centering
  \begin{tikzpicture}[x=1.15cm,y=0.88cm,>=Latex,font=\sffamily\small]
    \draw[vnink,line width=0.8pt,-{Latex[length=3mm]}] (0,0)--(7.2,0) node[right]{yield $r$};
    \draw[vnink,line width=0.8pt,-{Latex[length=3mm]}] (0,0)--(0,4.7) node[above]{price $\B(r)$};
    \draw[vnlink,line width=1.4pt,domain=0.45:6.7,samples=80,smooth]
      plot (\x,{5.25/(0.75+0.42*\x)});
    \coordinate (mid) at (3.45,{5.25/(0.75+0.42*3.45)});
    \coordinate (low) at (2.15,{5.25/(0.75+0.42*2.15)});
    \coordinate (high) at (4.75,{5.25/(0.75+0.42*4.75)});
    \fill[vnink] (mid) circle (1.8pt);
    \fill[vnlink] (low) circle (1.8pt);
    \fill[vncarnelian] (high) circle (1.8pt);
    \draw[vnrule,dashed] (2.15,0)--(low) (3.45,0)--(mid) (4.75,0)--(high);
    \node[below] at (2.15,0) {$r_-$};
    \node[below] at (3.45,0) {$r_0$};
    \node[below] at (4.75,0) {$r_+$};
    \draw[vnrule,dashed] (1.82,3.18)--(low) (1.82,2.39)--(mid);
    \draw[vnlink,<->,line width=0.9pt] (1.82,2.39)--(1.82,3.18)
      node[midway,left,text=vnlink]{gain};
    \draw[vnrule,dashed] (high)--(5.08,1.91) (mid)--(5.08,2.39);
    \draw[vncarnelian,<->,line width=0.9pt] (5.08,1.91)--(5.08,2.39)
      node[midway,right,text=vncarnelian]{loss};
    \node[anchor=west,text=vnmuted] at (5.05,3.35) {decreasing};
    \node[anchor=west,text=vnmuted] at (5.05,2.95) {and convex};
  \end{tikzpicture}
  \caption{The price-yield curve for a fixed positive cash-flow stream. The comparison yields are $r_-:=r_0-\Delta r$ and $r_+:=r_0+\Delta r$ for an equal change $\Delta r>0$ from the initial yield $r_0$. Convexity makes the vertical price gain from the rate decrease exceed the price loss from the equal rate increase.}
  \label{fig:price-yield}
\end{figure}

The bill-sensitivity notebook tests the first part of \thmref{yield-shape} using zero-coupon Treasury auction data, where the single maturity payment makes the mechanism especially transparent.

\notebooklink
  {L2b: Treasury Bill Price Sensitivity}
  {https://github.com/varnerlab/CHEME-5660-CourseRepository-Fall-2026/blob/main/lectures/week-2/L2b/CHEME-5660-L2b-Malkiel-TBill-Example-Fall-2025.ipynb}
  {Measure how a zero-coupon bill price changes with yield and remaining time using Treasury auction records.}

\section{Duration and Convexity}

The derivative signs answer \emph{which direction}; risk management also needs \emph{how much}. Normalizing the derivatives by price produces two local measures with interpretable units (\defnref{duration-convexity}).

\begin{definition}{Modified duration and convexity}{duration-convexity}
For the positive fixed cash flows, nominal annual yield $r$, coupon frequency $n$, and price $\B(r)>0$ defined in \defnref{yield-price}, the modified duration $D_{\mathrm{mod}}(r)$ and convexity $K(r)$ are defined by \eqnref{risk-measures}:
\begin{equation}
  \label{eq:risk-measures}
  D_{\mathrm{mod}}(r):=-\frac{1}{\B(r)}\frac{d\B}{dr},
  \qquad
  K(r):=\frac{1}{\B(r)}\frac{d^2\B}{dr^2}.
\end{equation}
When the yield $r$ is expressed as a decimal per year, modified duration has units of years and convexity has units of years squared. Both measures are positive under the assumptions of \thmref{yield-shape}.
\end{definition}

For a small yield shock $\Delta r$, a second-order Taylor approximation combines the two normalized sensitivities. The approximate fractional price change is given by \eqnref{duration-convexity-approximation}:
\begin{equation}
  \label{eq:duration-convexity-approximation}
  \frac{\Delta\B}{\B}
  \approx-D_{\mathrm{mod}}\Delta r+\frac{1}{2}K(\Delta r)^2.
\end{equation}
The duration term supplies the local slope, while the convexity term corrects the straight-line approximation for curvature. This formula uses a decimal yield change: a 25-basis-point shock is $\Delta r=0.0025$, where one basis point is $10^{-4}$.

\section{How Contract Structure Changes Sensitivity}

Duration is a present-value-weighted measure of payment timing. Moving more value toward distant dates increases exposure to discount-rate changes; moving value toward earlier dates reduces it. That principle gives the carefully qualified maturity and coupon comparisons associated with Malkiel.

\textbf{Maturity effect.}\enspace For zero-coupon securities, all value arrives at maturity, so a longer remaining maturity unambiguously increases duration and yield sensitivity. For otherwise comparable option-free coupon bonds, sensitivity generally increases with maturity, but at a decreasing rate. This is a comparison of yield sensitivity---not a claim that adding another coupon period always lowers the bond's price.

\textbf{Coupon effect.}\enspace Holding yield, maturity, face value, and coupon frequency fixed, a higher coupon shifts a larger fraction of present value toward earlier dates. The resulting percentage price sensitivity is therefore lower.

\textbf{Discrete-maturity caution.}\enspace A coupon count $N$ is an integer. Extending maturity adds a payment and changes the cash-flow vector, so differentiating a finite coupon sum with respect to a continuous maturity variable can be misleading. Compare two explicitly defined cash-flow streams or use duration with the cash flows held fixed.

The coupon-security notebook evaluates these comparative statics on Treasury note data and shows how coupon placement changes the response to a yield shock.

\newpage
\thispagestyle{fancy}
\notebooklink
  {L2b: Treasury Note and Bond Price Sensitivity}
  {https://github.com/varnerlab/CHEME-5660-CourseRepository-Fall-2026/blob/main/lectures/week-2/L2b/CHEME-5660-L2b-Malkiel-NotesBonds-Example-Fall-2025.ipynb}
  {Compare yield, coupon, and maturity sensitivities for coupon-bearing Treasury securities.}

\keytakeaways
  {In this lecture, we differentiated a fixed cash-flow bond price, established its decreasing convex price-yield shape, and converted those derivatives into duration and convexity risk measures.}
  {Yield and price oppose}
  {For a fixed positive cash-flow stream, a higher yield lowers price because every promised payment is discounted more heavily.}
  {Curvature creates asymmetry}
  {Positive convexity makes the gain from a yield decrease larger than the loss from an equal yield increase around the same starting yield.}
  {Timing controls exposure}
  {Longer-dated value increases interest-rate sensitivity, while larger coupons move value earlier and generally reduce percentage sensitivity.}
  {Next, use observed Treasury prices to infer yields across maturities and assemble the yield curve.}

\begin{readinglist}
  \item Burton G. Malkiel, \href{https://academic.oup.com/qje/article-abstract/76/2/197/1904062}{``Expectations, Bond Prices, and the Term Structure of Interest Rates,''} \emph{Quarterly Journal of Economics} 76(2), 197--218 (1962).
  \item Federal Reserve Board, \href{https://www.federalreserve.gov/publications/2023-April-SVB-Key-Takeaways.htm}{\emph{Review of the Federal Reserve's Supervision and Regulation of Silicon Valley Bank}} (2023), for the interaction of interest-rate, liquidity, and governance risk.
  \item Frank J. Fabozzi and Francesco A. Fabozzi, \href{https://mitpress.mit.edu/9780262046275/bond-markets-analysis-and-strategies/}{\emph{Bond Markets, Analysis, and Strategies}}, 10th ed. (2021), for duration and convexity conventions.
\end{readinglist}

\vfill
{\sffamily\footnotesize\color{vnmuted}\textbf{Educational use and risk notice.} These notes are for instruction only and do not constitute investment advice. Financial decisions involve risk, and model outputs depend on assumptions.}

\end{document}
